% CN 8.10 Exercises (p.299--304)

% ======================================================================
\begin{exercise}{1}
    Consider the following algorithm.
    \begin{quote}
        for $i$ from 1 to $a$ do\\
        \hspace*{1em} for $j$ from 1 to $b$ do\\
        \hspace*{2em} for $k$ from 1 to $c$ do\\
        \hspace*{3em} for $l$ from 1 to $d$ do\\
        \hspace*{4em} beep!\\
        \hspace*{1em} for $m$ from 1 to $e$ do\\
        \hspace*{2em} for $n$ from 1 to $f$ do\\
        \hspace*{3em} beep!\\
        \hspace*{2em} for $o$ from 1 to $g$ do\\
        \hspace*{3em} beep!\\
        \hspace*{2em} for $p$ from 1 to $h$ do\\
        \hspace*{3em} beep!
    \end{quote}
    \begin{enumerate}[label=(\alph*)]
        \item Give an expression in $a,b,c,d,e,f,g,h$ for the number of times
              this algorithm beeps.
        \item If each of these eight variables is $O(n)$, give a big-O expression
              for the number of beeps.
    \end{enumerate}
\end{exercise}

\begin{exercise}{2}
    \begin{enumerate}[label=(\alph*)]
        \item A House of Representatives ballot paper has $n$ boxes, one per
              candidate. A voter must enter the numbers $1$ to $n$ in order of
              preference, exactly one number per box. In how many ways can this
              be done?
        \item A Senate ballot paper has one box for each of $N$ candidates, but
              voters need only enter the numbers $1$ to $12$ for their twelve most
              preferred candidates. In how many ways can this be done?
    \end{enumerate}
\end{exercise}

\begin{exercise}{3}
    Two fair dice are thrown; the outcome is the ordered pair of numbers shown.
    \begin{enumerate}[label=(\alph*)]
        \item How many possible outcomes are there?
        \item How many outcomes have the first number less than the second?
        \item How many outcomes have sum 7?
        \item How many outcomes have both numbers $\le3$?
        \item How many outcomes have neither number equal to 5?
    \end{enumerate}
\end{exercise}

% ======================================================================
\begin{exercise}{4}
    The Melbourne city centre grid has five ``horizontal'' streets (Flinders,
    Collins, Bourke, Lonsdale, La Trobe, from south to north) and nine
    ``vertical'' streets (Spencer, King, William, Queen, Elizabeth, Swanston,
    Russell, Exhibition, Spring, from west to east), giving $4\times8=32$ square
    blocks.

    You start at the corner of Flinders and Swanston, and walk along main streets,
    never turning back, always going further away from your starting point.
    \begin{enumerate}[label=(\alph*)]
        \item How many routes are there to the corner of William and La Trobe?
              Express your answer using a counting expression, then evaluate it.
        \item How many such routes do not meet Russell Street and take you four
              block-lengths away? Give the total and the number of routes to each
              intersection at this distance. Comment on the relationship among
              these numbers.
        \item How many shortest routes are there from the south-west corner
              (Flinders and Spencer) to the north-east corner (Spring and La
              Trobe)?
    \end{enumerate}
\end{exercise}

\begin{exercise}{5}
    Using only Basic Searches and the inclusion-exclusion principle, determine how
    many Monash Library catalogue entries contain at least one of the three terms
    CSIRAC, SILLIAC and WREDAC.
\end{exercise}

\begin{exercise}{6}
    \begin{enumerate}[label=(\alph*)]
        \item Use inclusion-exclusion to determine how many positive integers
              $\le100$ are not multiples of 3, 5 or 7.
        \item Is there another simple way of working this out, to check your
              answer?
    \end{enumerate}
\end{exercise}

\begin{exercise}{7}
    Use inclusion-exclusion to determine the number of surjections from
    $\{a,b,c,d,e,f\}$ to $\{1,2,3,4\}$.
\end{exercise}

\begin{exercise}{8}
    Determine the numbers of $n$-digit positive integers that have their digits
    \begin{enumerate}[label=(\alph*)]
        \item in ascending order;
        \item in descending order;
        \item in nondecreasing order;
        \item in nonincreasing order.
    \end{enumerate}
\end{exercise}

% ======================================================================
\begin{exercise}{9}
    \textbf{Poker.} A standard deck has 52 cards in four suits
    $\spadesuit,\clubsuit,\diamondsuit,\heartsuit$ of 13 cards each, with ranks
    $2,3,\dots,10,J,Q,K,A$ ($J,Q,K,A=11,12,13,14$). The Ace can also stand for 1,
    but not both low and high in the same hand. A poker hand is a set of five
    cards. Five cards are \emph{consecutive} if their ranks are in numerical
    sequence (A,2,3,4,5 and 10,J,Q,K,A are allowed, but there is no wrap-around,
    e.g.\ Q,K,A,2,3 is not consecutive).

    How many poker hands are there of each type:
    \begin{enumerate}[label=(\alph*)]
        \item no restriction at all;
        \item straight flush: five consecutive cards in the same suit;
        \item four of a kind: four of the five cards have the same rank;
        \item full house: three cards of one rank, two of another;
        \item flush: all five cards the same suit, ranks not all consecutive;
        \item straight: five consecutive cards, not all in the same suit;
        \item three of a kind: three of the same rank, the other two of
              different ranks to those three and to each other;
        \item pair of pairs: two pairs of different ranks, and a fifth card of a
              different rank to the others;
        \item one pair: a single pair, the other three cards of different ranks
              to the pair and to each other;
        \item nothing: none of the types (b)--(i).
    \end{enumerate}
\end{exercise}

% ======================================================================
\begin{exercise}{10}
    How many strings of five digits are there in which
    \begin{enumerate}[label=(\alph*)]
        \item there is no restriction whatsoever;
        \item leading 0s are not allowed;
        \item all digits are identical;
        \item all digits are different;
        \item every digit is different to its predecessor (other repetitions
              allowed);
        \item all digits are in strictly ascending order;
        \item all digits are in nondecreasing order;
        \item the ordering of digits is monotonic (either nondecreasing or
              nonincreasing);
        \item the digits are in arithmetic progression;
        \item the digits are in geometric progression;
        \item the digits add up to a multiple of 9;
        \item no digit appears in its own position (the first digit is not 1,
              the second is not 2, and so on);
        \item in addition to (l), all the digits are different.
    \end{enumerate}
\end{exercise}

\begin{exercise}{11}
    Specify formally the bijection given in \S8.9
    \begin{itemize}
        \item from the set of $x$-sequences of nonnegative integers satisfying
              the constraints given there (including (8.12))
        \item to the set of $y$-sequences of positive integers satisfying the
              constraints given there (including (8.13)).
    \end{itemize}
\end{exercise}
